\subsection{Intersectional Subgroup Performance}
\label{sec:subgroup_performance}

\begin{table*}[htbp]
\caption{Subgroup-Level Performance for Tri-Domain Models.}
\label{tab:XI}
\centering
\begin{tabular}{llcccc}
\toprule
Classifier & Subgroup & Recall & Precision & F1-Score & OvR Accuracy \\
\midrule
\multirow{6}{*}{\textbf{SVM (Tri)}} & White Males & 96.94\% & 95.36\% & 96.14\% & 98.70\% \\
 & Black Males & 94.17\% & 95.49\% & 94.83\% & 98.29\% \\
 & White Females & 94.72\% & 92.41\% & 93.55\% & 97.82\% \\
 & Asian Males & 94.44\% & 92.39\% & 93.41\% & 97.78\% \\
 & Asian Females & 92.50\% & 92.50\% & 92.50\% & 97.50\% \\
 & Black Females & 89.44\% & 94.15\% & 91.74\% & 97.31\% \\
\midrule
\multirow{6}{*}{\textbf{LR (Tri)}} & White Males & 96.11\% & 95.05\% & 95.58\% & 98.52\% \\
 & Black Males & 93.06\% & 95.71\% & 94.37\% & 98.15\% \\
 & White Females & 92.22\% & 91.21\% & 91.71\% & 97.22\% \\
 & Asian Males & 92.78\% & 90.76\% & 91.76\% & 97.22\% \\
 & Asian Females & 91.11\% & 91.62\% & 91.36\% & 97.13\% \\
 & Black Females & 91.11\% & 92.13\% & 91.62\% & 97.22\% \\
\bottomrule
\end{tabular}
\end{table*}

Profil kinerja klasifikasi granular pada model terbaik SVM berbasis fusi tri-domain Face $\oplus$ Emotion $\oplus$ Age dievaluasi secara mendalam berdasarkan metrik OvR sebagaimana disajikan pada Table~\ref{tab:XI}. Seluruh subkelompok demografis berhasil mencapai F1-Score melampaui 91.00\%, dengan rentang capaian bergerak dari 91.74\% pada Black Females hingga 96.14\% pada White Males, sehingga menghasilkan selisih disparitas rentang sebesar 4.40\%. Sementara itu, nilai Akurasi OvR berada pada rentang 97.31\% hingga 98.70\% dengan disparitas sebesar 1.39 pp. Perlu dicatat secara metodologis bahwa tingginya nilai Akurasi OvR tersebut turut dipengaruhi oleh rasio dominansi sampel negatif sebesar 5:1 pada evaluasi biner OvR, sehingga metrik ini tidak diposisikan sebagai indikator tunggal disparitas kinerja.

Perbandingan profil kinerja subkelompok antara model SVM dan pengklasifikasi linier LR pada konfigurasi tri-domain disajikan pada Table~\ref{tab:XI}. Model LR tri-domain turut memperlihatkan konsistensi performa yang tinggi dengan capaian F1-Score melampaui 91.00\% pada seluruh kelas demografis, yang merentang dari 91.36\% pada Asian Females hingga 95.58\% pada White Males dengan selisih disparitas rentang sebesar 4.22\%. Subkelompok White Males secara konsisten mencatatkan performa klasifikasi tertinggi pada kedua pengklasifikasi tersebut. Temuan empiris ini membuktikan bahwa fusi representasi tri-domain menghasilkan representasi visual yang tangguh dan terdistribusi stabil melintasi seluruh subkelompok demografis, baik pada model linier maupun model non-linier. Meskipun demikian, evaluasi subkelompok ini mencerminkan kinerja granular per kelas dan bukan audit keadilan algoritmik komprehensif.
